\documentclass[12pt,a4paper]{article}

\usepackage[utf8]{inputenc}
\usepackage[T1]{fontenc}
\usepackage[turkish]{babel}
\usepackage{geometry}
\usepackage{hyperref}
\usepackage{listings}
\usepackage{xcolor}
\usepackage{booktabs}
\usepackage{tabularx}
\usepackage{enumitem}
\usepackage{titlesec}
\usepackage{fancyhdr}
\usepackage{graphicx}
\usepackage{microtype}
\usepackage{parskip}

\geometry{
  top=2.5cm, bottom=2.5cm,
  left=2.8cm, right=2.8cm
}

% Kod bloğu stili
\definecolor{codebg}{HTML}{F5F5F5}
\definecolor{codeborder}{HTML}{CCCCCC}
\definecolor{keyword}{HTML}{0000FF}
\definecolor{comment}{HTML}{6A9955}
\definecolor{string}{HTML}{A31515}

\lstset{
  backgroundcolor=\color{codebg},
  basicstyle=\ttfamily\small,
  breaklines=true,
  frame=single,
  rulecolor=\color{codeborder},
  keywordstyle=\color{keyword},
  commentstyle=\color{comment},
  stringstyle=\color{string},
  showstringspaces=false,
  xleftmargin=4pt,
  xrightmargin=4pt,
}

% Başlık formatı
\titleformat{\section}{\large\bfseries}{}{0em}{\thesection\quad}[\titlerule]
\titleformat{\subsection}{\normalsize\bfseries}{}{0em}{\thesubsection\quad}

% Header / footer
\pagestyle{fancy}
\fancyhf{}
\rhead{\small DevOps Case Study — Teknik Rapor}
\lhead{\small Deniz YALDIZ}
\cfoot{\thepage}

% Hyperlink renkleri
\hypersetup{
  colorlinks=true,
  linkcolor=black,
  urlcolor=blue,
}

% ─────────────────────────────────────────────────────────────
\begin{document}

\begin{titlepage}
  \centering
  \vspace*{2cm}
  {\Huge\bfseries DevOps Case Study\par}
  \vspace{0.5cm}
  {\Large Teknik Uygulama Raporu\par}
  \vspace{2cm}
  {\large\bfseries Deniz YALDIZ\par}
  \vspace{0.3cm}
  {\normalsize Mayıs 2026\par}
  \vspace{3cm}
  \begin{tabular}{ll}
    \textbf{Platform}   & Linux (Ubuntu 22.04) \\
    \textbf{K8s}        & v1.29.15 \\
    \textbf{Araçlar}    & Docker, Vagrant, GitLab CI, Prometheus, Grafana \\
    \textbf{Repo}       & \texttt{github.com/DenizYald3iz/devops-case-study} \\
  \end{tabular}
  \vfill
  {\small Bu rapor her sorudaki teknik kararların gerekçesini,
  değerlendirilen alternatifleri ve trade-off analizini içermektedir.}
\end{titlepage}

\tableofcontents
\newpage

% ─────────────────────────────────────────────────────────────
\section{Genel Mimari}

Proje beş bağımsız bileşenden oluşmaktadır. Her bileşen kendi
klasöründe izole edilmiş ve bağımsız olarak çalıştırılabilecek
şekilde tasarlanmıştır.

\begin{center}
\begin{tabular}{lll}
\toprule
\textbf{Bileşen} & \textbf{Teknoloji} & \textbf{Çalışma Ortamı} \\
\midrule
Uygulama container'ı & Docker + Streamlit 1.32 & Docker Compose \\
K8s cluster           & Vagrant + VirtualBox     & 3 VM (1 master, 2 worker) \\
K8s deployment        & kubectl manifests        & case-study namespace \\
CI/CD pipeline        & GitLab CI                & GitLab Runner \\
Monitoring            & Prometheus + Grafana     & Docker Compose \\
\bottomrule
\end{tabular}
\end{center}

Genel trafik akışı:

\begin{lstlisting}
git push (main)
    │
    ▼
GitLab CI  →  Kaniko build  →  Trivy scan  →  Registry push
                                                     │
                                              kubectl set image
                                                     │
                                                     ▼
Kullanıcı :80  →  NodePort  →  Service:8080  →  Pod:8501 (Streamlit)
                                                     │
                                              cAdvisor / node-exporter
                                                     │
                                              Prometheus  →  Grafana
\end{lstlisting}

% ─────────────────────────────────────────────────────────────
\section{Soru 1 — Containerization}

\subsection{Python Versiyon Seçimi}

Şablon dosyasında \texttt{python:3.8-slim} kullanılıyordu. Python 3.8,
Ekim 2024 itibarıyla End-of-Life olmuştur; güvenlik yamaları artık
yayınlanmamaktadır. Bu nedenle \texttt{python:3.12-slim}'e geçildi.
Tercih raporda açıkça belirtilmiştir.

\subsection{Multi-Stage Build}

\begin{lstlisting}[language=docker,caption={Multi-stage build yapısı}]
FROM python:3.12-slim AS builder
WORKDIR /app
COPY requirements.txt .
RUN pip install --user --no-cache-dir -r requirements.txt

FROM python:3.12-slim
# Yalnizca bağimliliklar kopyalanir, pip ve build araclari gelmez
COPY --from=builder /root/.local /home/appuser/.local
\end{lstlisting}

Build araçları ve pip final image'a taşınmadığı için image boyutu
anlamlı ölçüde küçülmektedir.

\subsection{Non-Root Kullanıcı ve UID Sabitleme}

Güvenlik temeli olarak container root yerine \texttt{appuser} ile
çalışmaktadır. Kritik bir detay: \texttt{useradd -r} ile otomatik
atanan UID, base image'daki mevcut sistem kullanıcılarına bağlı
olarak değişebilir. Kubernetes deployment YAML'ında
\texttt{runAsUser: 999} yazıldığı için bu değerin container ile
örtüşmesi zorunludur. Bu nedenle UID ve GID açıkça sabitlendi:

\begin{lstlisting}[language=bash]
RUN groupadd -r -g 999 appuser && \
    useradd -r -u 999 -g appuser -m appuser
\end{lstlisting}

\subsection{ENTRYPOINT ve Başlangıç Logu}

\texttt{CMD} yerine ayrı bir \texttt{entrypoint.sh} kullanılmaktadır.
Bu tercih iki avantaj sağlar:

\begin{enumerate}[leftmargin=*]
  \item Konteyner başladığında hangi portta çalıştığını log'a yazdırır.
  \item \texttt{STREAMLIT\_SERVER\_PORT} environment variable'ını
        \texttt{-{}-server.port} parametresine iletir; port konfigürasyonu
        tek noktadan yönetilir.
\end{enumerate}

\subsection{Port Konfigürasyonu (.env Entegrasyonu)}

Başlangıçta \texttt{docker-compose.yml} içinde port değerleri
hardcoded yazılıyordu (\texttt{"8501:8501"}). Bu durum
\texttt{.env} dosyasındaki \texttt{STREAMLIT\_SERVER\_PORT}
değerini değiştirmenin etkisiz kalmasına yol açıyordu.
Düzeltme ile birlikte:

\begin{lstlisting}[language=yaml]
ports:
  - "${STREAMLIT_SERVER_PORT:-8501}:${STREAMLIT_SERVER_PORT:-8501}"
environment:
  - STREAMLIT_SERVER_PORT=${STREAMLIT_SERVER_PORT:-8501}
\end{lstlisting}

\subsection{Healthcheck}

Dockerfile'da \texttt{HEALTHCHECK} ve Compose'da \texttt{healthcheck}
bloğu tanımlıdır. Her ikisi de Streamlit'in built-in
\texttt{/\_stcore/health} endpoint'ini kullanmaktadır.

% ─────────────────────────────────────────────────────────────
\section{Soru 2 — Kubernetes Cluster (Vagrant)}

\subsection{Donanım ve Distro Seçimi}

\begin{center}
\begin{tabular}{llll}
\toprule
\textbf{Node} & \textbf{IP} & \textbf{CPU} & \textbf{RAM} \\
\midrule
k8s-master   & 192.168.56.10 & 2 vCPU & 2 GB \\
k8s-worker-1 & 192.168.56.11 & 2 vCPU & 2 GB \\
k8s-worker-2 & 192.168.56.12 & 2 vCPU & 2 GB \\
\bottomrule
\end{tabular}
\end{center}

Ubuntu 22.04 LTS (\texttt{ubuntu/jammy64}) tercih edilmiştir.
kubeadm referans dokümantasyonunda en geniş test edilen dağıtımdır.
Toplam kaynak ihtiyacı yaklaşık 6 GB RAM ve 6 vCPU'dur.

\subsection{Container Runtime: containerd}

Docker değil, containerd kullanılmaktadır. Kubernetes 1.24 sürümüyle
birlikte dockershim kaldırılmıştır; containerd doğrudan CRI uyumludur.
\texttt{SystemdCgroup = true} ayarı zorunludur — aksi takdirde kubelet
ve containerd farklı cgroup driver'ı kullandığında kaynak yönetimi
çakışmalarına yol açar.

\subsection{CNI: Calico}

Flannel ile karşılaştırıldığında Calico'nun tercih edilme gerekçesi
\texttt{NetworkPolicy} desteğidir. Pod-to-pod trafik kısıtlaması
üretim ortamlarında standart bir gereklilik olduğundan,
NetworkPolicy desteği olmayan Flannel elendi.

\subsection{Karşılaşılan Sorunlar ve Çözümler}

\textbf{GPG /dev/tty hatası:}
\texttt{gpg -{}-dearmor} komutu, Vagrant'ın non-interactive SSH
ortamında \texttt{/dev/tty} açmaya çalışmaktadır. Provision
bağlantısında bu cihaz mevcut olmadığından \texttt{curl: (23) Failed
writing body} hatası zincirleme olarak tetiklenmektedir. Çözüm olarak
\texttt{-{}-batch -{}-yes -{}-no-tty} parametreleri eklendi.

\textbf{vagrant provision idempotency:}
\texttt{kubeadm init}, ikinci çalıştırmada portların ve manifest
dosyalarının mevcut olduğunu görüp hata vermektedir. Çözüm olarak
\texttt{/etc/kubernetes/admin.conf} varlığı kontrol edilerek init adımı
koşullu hale getirildi. Worker'larda aynı şekilde
\texttt{/etc/kubernetes/kubelet.conf} kontrolü eklendi.

\subsection{Sabit Token ile Join (Idempotency)}

\texttt{kubeadm-config.yaml} içinde önceden tanımlı token kullanılmaktadır.
Bu sayede \texttt{master.sh} tamamlandıktan sonra join komutu öngörülebilir
formatta \texttt{/vagrant/join.sh}'a yazılabilmekte, worker'lar bu
dosyayı bekleyerek cluster'a katılmaktadır.

% ─────────────────────────────────────────────────────────────
\section{Soru 3 — Kubernetes Deployment}

\subsection{Port 80 Problemi ve Çözüm Alternatifleri}

NodePort'un varsayılan aralığı 30000--32767'dir; 80 numaralı port bu
aralığın dışındadır. Üç çözüm değerlendirilmiştir:

\begin{center}
\begin{tabularx}{\textwidth}{lXll}
\toprule
\textbf{Yöntem} & \textbf{Açıklama} & \textbf{vagrant up sonrası ek adım} & \textbf{Üretim uygunluğu} \\
\midrule
A — apiserver range & \texttt{-{}-service-node-port-range=80-32767} & Yok & Düşük \\
B — Ingress Controller & NGINX, host/path bazlı routing & Var (controller kurulumu) & Yüksek \\
C — iptables NAT & Host'ta 80$\to$30080 yönlendirme & Var, kalıcı değil & Düşük \\
\bottomrule
\end{tabularx}
\end{center}

\textbf{Seçilen yöntem: A.}
Demo ortamı için temel kriter şuydu: \texttt{vagrant up} tamamlandıktan
sonra ek adım olmaksızın \texttt{curl http://192.168.56.10/} çalışsın.
Yöntem A bu kriteri karşılamaktadır.

\textbf{Trade-off:} Bu ayar cluster genelinde geçerlidir. Başka bir
servis port 80'i yanlışlıkla alabilir. Tek kiracı, kontrollü bir demo
ortamında kabul edilebilir; çok ekipli üretim ortamında kullanılmaz.

\texttt{ingress.yaml} dosyası hazır olarak repoya eklenmiştir. Üretim
ortamına geçişte yalnızca NGINX Ingress Controller kurulumu yeterlidir.

\subsection{Deployment Kararları}

\textbf{Rolling update stratejisi:}
\texttt{maxUnavailable: 0} ve \texttt{maxSurge: 1} kombinasyonuyla
rolling update sırasında her zaman 2 replica ayakta kalmaktadır.

\textbf{Probe stratejisi:}
\texttt{liveness} probe 30 saniye gecikmeyle başlamaktadır. Streamlit'in
tam açılması bu süreyi alabilmektedir; erken probe container'ı gereksiz
yere yeniden başlatır.

\textbf{readOnlyRootFilesystem: false:}
Streamlit \texttt{/tmp} altına geçici dosyalar yazmaktadır. Bu değer
\texttt{true} yapıldığında uygulama başlatılamamaktadır. Tüm
yetenekler \texttt{capabilities: drop: ["ALL"]} ile kaldırılarak
telafi edilmiştir.

\textbf{UID hizalaması:}
Deployment YAML'ında \texttt{runAsUser: 999} yazılmaktadır. Dockerfile'da
UID ve GID açıkça 999 olarak sabitlenmiştir; böylece Kubernetes güvenlik
bağlamı ile container kullanıcısı örtüşmektedir.

% ─────────────────────────────────────────────────────────────
\section{Soru 4 — GitLab CI/CD Pipeline}

\subsection{Pipeline Yapısı}

6 stage'den oluşan bir pipeline tasarlanmıştır:

\begin{lstlisting}
validate -> build -> scan -> push -> deploy -> rollback
\end{lstlisting}

\begin{center}
\begin{tabularx}{\textwidth}{llX}
\toprule
\textbf{Stage} & \textbf{Job} & \textbf{Açıklama} \\
\midrule
validate  & validate\_dockerfile  & Hadolint ile Dockerfile lint \\
build     & build\_image\_kaniko  & Kaniko ile rootless image build + SHA tag push \\
scan      & security\_scan        & Trivy HIGH/CRITICAL zafiyet taraması \\
push      & tag\_latest\_main     & SHA $\to$ latest tag (yalnızca main branch) \\
deploy    & deploy\_production    & kubectl rolling update — \textbf{manuel onay} \\
rollback  & rollback\_production  & kubectl rollout undo — \textbf{manuel, acil durum} \\
\bottomrule
\end{tabularx}
\end{center}

\subsection{Kaniko (DinD Yerine)}

Docker-in-Docker, runner'ın \texttt{privileged: true} modunda
çalışmasını zorunlu kılar; bu durum host kernel'e tam erişim anlamına
gelir. Kaniko aynı işi userspace'de gerçekleştirmekte, root yetkisi
veya özel socket erişimi gerektirmemektedir.

\subsection{Shift-Left Security: Trivy}

Güvenlik açığı bulunan bir image asla \texttt{latest} tag almaz ve
üretime gidemez. Pipeline sırası bunu zorunlu kılar:
\texttt{build $\to$ scan $\to$ push(latest) $\to$ deploy.}

\texttt{HIGH} veya \texttt{CRITICAL} zafiyet tespit edildiğinde
\texttt{-{}-exit-code 1} ile pipeline durdurulmaktadır.

\subsection{Continuous Delivery (Otomatik Deploy Değil)}

Deploy aşaması manuel onay gerektirmektedir. Üretim ortamına çıkış
her zaman bir insanın onayına bağlı kalmalıdır. Rollback stage,
terminal açmadan GitLab arayüzü üzerinden tek tıkla
\texttt{kubectl rollout undo} çalıştırma imkânı sağlamaktadır.

\subsection{Secret Yönetimi}

\texttt{KUBECONFIG\_PROD} değişkeni GitLab'da Masked ve Protected
olarak saklanmaktadır. Pipeline içinde geçici dosyaya decode
edilmekte, job tamamlandığında silinmektedir; log'a yazdırılmamaktadır.

\subsection{Düzeltilen Bug: Dockerfile Yolu}

İlk versiyonda \texttt{-{}-context \$\{CI\_PROJECT\_DIR\}} ve
\texttt{-{}-dockerfile \$\{CI\_PROJECT\_DIR\}/Dockerfile} yazılıydı.
Dockerfile proje kökünde değil \texttt{01-docker/} altında
bulunduğundan pipeline build aşamasında hata veriyordu.
Doğru yollar: \texttt{01-docker/} context ve
\texttt{01-docker/Dockerfile}.

% ─────────────────────────────────────────────────────────────
\section{Soru 5 — Monitoring}

\subsection{Araç Seçimi: Prometheus + Grafana}

Üç alternatif değerlendirilmiştir:

\begin{itemize}[leftmargin=*]
  \item \textbf{Custom shell script:} Sürdürülebilir değil, ekip
        tarafından kolayca genişletilemiyor, alert mekanizması yok.
  \item \textbf{Zabbix:} Ek kurulum karmaşıklığı yüksek; case'in
        odağı değil.
  \item \textbf{Prometheus + Grafana:} Endüstri standardı, container
        ve host metrikleri için olgun ekosistem, provisioning ile
        sıfır manuel kurulum.
\end{itemize}

\subsection{Bileşenler}

\begin{center}
\begin{tabular}{llll}
\toprule
\textbf{Araç} & \textbf{Versiyon} & \textbf{Görev} & \textbf{Port} \\
\midrule
Prometheus     & 2.51.0 & Metrik toplama, TSDB, alert değerlendirme & 9090 \\
Grafana        & 10.4.0 & Görselleştirme, provisioning              & 3000 \\
cAdvisor       & 0.49.1 & Container CPU/memory metrikleri           & 8080 \\
node-exporter  & 1.7.0  & Host CPU/memory/disk metrikleri           & 9100 \\
\bottomrule
\end{tabular}
\end{center}

\subsection{Dashboard}

"Streamlit App — CPU \& Memory" başlıklı dashboard altı panel
içermektedir: Container CPU \%, Container Memory (kullanım + limit),
Host CPU gauge, Host Memory Available gauge, Container Last Start
zamanı, Host Disk kullanımı.

Dashboard JSON, \texttt{grafana/dashboards/} altında saklanmakta
ve Grafana başladığında provisioning sistemi tarafından otomatik
yüklenmektedir. Manuel import gerekmemektedir.

\subsection{Alert Kuralları}

\begin{center}
\begin{tabular}{llll}
\toprule
\textbf{Alert} & \textbf{Eşik} & \textbf{Süre} & \textbf{Seviye} \\
\midrule
ContainerHighCPU    & $>80\%$         & 2 dakika & warning  \\
ContainerHighMemory & limit'in $>80\%$& 2 dakika & warning  \\
ContainerDown       & görünmüyor      & 1 dakika & critical \\
HostHighCPU         & $>85\%$         & 5 dakika & warning  \\
HostLowMemory       & $<15\%$ available & 5 dakika & critical \\
\bottomrule
\end{tabular}
\end{center}

\subsection{Karşılaşılan Sorunlar ve Çözümler}

\textbf{host.docker.internal (Linux):}
Docker Desktop'ın (Mac/Windows) otomatik çözdüğü bu hostname, Linux'ta
varsayılan olarak kayıtlı değildir. Prometheus servisine
\texttt{extra\_hosts: ["host.docker.internal:host-gateway"]} eklenerek
çözüldü.

\textbf{Sahte Streamlit metrics endpoint'i:}
\texttt{/\_stcore/metrics} Streamlit'te bulunmamaktadır; bu job her
scrape işleminde 404 alıyordu. Job kaldırıldı. Container metrikleri
cAdvisor üzerinden \texttt{name="streamlit-app"} label'ı ile
sorgulanmaktadır.

\textbf{1970 tarihi (Container Last Start paneli):}
\texttt{container\_start\_time\_seconds} saniye cinsinden döner;
Grafana'nın \texttt{dateTimeAsIso} birimi milisaniye bekler.
\texttt{* 1000} çarpımıyla düzeltildi.

% ─────────────────────────────────────────────────────────────
\section{Üretim Ortamına Geçişte Ek Adımlar}

Case kapsamı dışında kalan ancak üretim ortamında gerekli olacak
bileşenler:

\begin{itemize}[leftmargin=*]
  \item \textbf{Secret yönetimi:} HashiCorp Vault veya Sealed Secrets.
        \texttt{KUBECONFIG} gibi değerler GitLab variable yerine
        merkezi secret store'dan alınmalıdır.
  \item \textbf{Log aggregation:} Loki + Promtail. Prometheus metrik,
        Loki log; her ikisi Grafana'da tek panelde gösterilebilir.
  \item \textbf{Helm chart:} Deployment manifest'leri Helm chart'a
        taşınarak ortam bazlı \texttt{values.yaml} ile parametrize
        edilebilir.
  \item \textbf{HorizontalPodAutoscaler:} CPU/memory eşiğine göre
        replica sayısını otomatik artırır.
  \item \textbf{Pod Disruption Budget:} Node bakımı sırasında minimum
        replica garantisi.
  \item \textbf{Network Policy:} Calico kullanıldığından ek altyapı
        gerektirmeksizin pod-to-pod trafik kısıtlaması uygulanabilir.
  \item \textbf{Alertmanager:} Şu an kurallar tanımlı ancak
        Alertmanager servisine yönlendirilmiyor. Slack/PagerDuty
        entegrasyonu için gerekli.
  \item \textbf{TLS:} Ingress üzerinde cert-manager + Let's Encrypt
        ile otomatik sertifika yönetimi.
\end{itemize}

% ─────────────────────────────────────────────────────────────
\section{Test Sonuçları}

\subsection{Soru 1}
\begin{lstlisting}
$ docker compose up -d --build
$ curl -o /dev/null -w "%{http_code}" http://localhost:8501/_stcore/health
200
\end{lstlisting}

\subsection{Soru 2}
\begin{lstlisting}
$ vagrant up
$ vagrant ssh k8s-master -c "kubectl get nodes"
NAME           STATUS   ROLES           AGE   VERSION
k8s-master     Ready    control-plane   17m   v1.29.15
k8s-worker-1   Ready    <none>          16m   v1.29.15
k8s-worker-2   Ready    <none>          14m   v1.29.15
\end{lstlisting}

\subsection{Soru 3}
\begin{lstlisting}
$ kubectl get deployment streamlit-app -n case-study
NAME            READY   UP-TO-DATE   AVAILABLE   AGE
streamlit-app   2/2     2            2           2m

$ kubectl get service streamlit-app -n case-study
NAME            TYPE       CLUSTER-IP     PORT(S)       AGE
streamlit-app   NodePort   10.107.88.91   8080:80/TCP   2m

$ curl -o /dev/null -w "%{http_code}" http://192.168.56.10/_stcore/health
200
$ curl -o /dev/null -w "%{http_code}" http://192.168.56.11/_stcore/health
200
$ curl -o /dev/null -w "%{http_code}" http://192.168.56.12/_stcore/health
200

# Güvenlik doğrulama
$ kubectl exec -n case-study <pod> -- id
uid=999(appuser) gid=999(appuser) groups=999(appuser)
\end{lstlisting}

\subsection{Soru 5}
\begin{lstlisting}
# Prometheus targets
Job: cadvisor        State: up
Job: node-exporter   State: up
Job: prometheus      State: up

# Uygulama metrikleri
streamlit-app memory : 70.7 MB
streamlit-app CPU    : 0.000% (idle)
Host CPU usage       : 16.0%
Host memory available: 12.5%
Host disk usage      : 77.3%

# Alert durumu
ContainerHighCPU    : inactive
ContainerHighMemory : inactive
ContainerDown       : inactive
HostHighCPU         : inactive
HostLowMemory       : pending  (12.5% < 15% eşiği)
\end{lstlisting}

% ─────────────────────────────────────────────────────────────
\section{Karşılaşılan Zorluklar}

\begin{enumerate}[leftmargin=*]
  \item \textbf{GPG provision hatası (Soru 2):} Vagrant SSH
        bağlantısının non-interactive olduğu göz önüne alınmadan
        yazılmış GPG komutu \texttt{/dev/tty} açmaya çalışıyordu.
        Hata mesajı doğrudan GPG'ye değil, zincirleme olarak curl'e
        işaret ettiğinden root cause bulmak zaman aldı.

  \item \textbf{Idempotency sorunu (Soru 2):} \texttt{vagrant provision}
        çalıştırıldığında kubeadm'ın halihazırda kurulu cluster'ı
        görüp hata vermesi beklenmedik bir davranıştı. Çözüm
        \texttt{admin.conf} varlık kontrolüyle sağlandı.

  \item \textbf{Dockerfile yolu (Soru 4):} CI pipeline'ındaki
        \texttt{-{}-context} ve \texttt{-{}-dockerfile} yolları proje
        kökünü gösteriyordu; Dockerfile \texttt{01-docker/} altında
        olduğundan build aşaması başarısız oluyordu. Çalıştırılıp
        hata görülmeden fark edilmesi zor bir bug.

  \item \textbf{host.docker.internal (Soru 5):} Docker Desktop'ta
        çalışan örnekler baz alınarak yazılmış Prometheus konfigürasyonu
        Linux'ta doğrudan çalışmadı. \texttt{extra\_hosts:
        host-gateway} çözümü Linux-spesifik bir detay.
\end{enumerate}

% ─────────────────────────────────────────────────────────────
\section{Sonuç}

Beş sorunun tamamı çalışır halde teslim edilmektedir. Her kararın
ardında bir gerekçe ve değerlendirilen alternatifler mevcuttur.
Demo ortamında bazı trade-off'lar kasıtlı olarak kolaylık yönünde
yapıldı (NodePort range genişletme gibi); bunların üretim ortamında
neden farklı yapılması gerektiği ilgili README dosyalarında açıklandı.

Çözümlerin tamamı tekrarlanabilir şekilde tasarlandı: \texttt{vagrant up}
tek komutla üç node'lu cluster'ı ayağa kaldırmakta,
\texttt{docker compose up} monitoring stack'ini başlatmakta,
\texttt{kubectl apply} deployment'ı oluşturmaktadır.

\end{document}
